% =====================================================================
%  minggu-06.tex  --  Slide Pertemuan Minggu ke-6
%  Pengantar Kriptografi (MPFA014) -- Prodi Matematika FMIPA UBB
%
%  Topik silabus : One-time pad
%  Sumber materi : buku_ajar/chapters/04/section-06.tex (Bab 4, \S6)
%  Sub-CPMK      : 1.2
%
%  Kompilasi dari dalam folder ppt/:
%      pdflatex -interaction=nonstopmode minggu-06.tex   (dua kali)
% =====================================================================

\documentclass[aspectratio=169,11pt]{beamer}
% =====================================================================
%  beamer-preamble.tex  --  prasarana dan gaya slide kuliah
%
%  Mata kuliah : Pengantar Kriptografi (MPFA014)
%  Program     : Program Studi Matematika, FMIPA, Universitas Bale Bandung
%
%  Berkas ini TIDAK memuat \documentclass. Setiap berkas minggu memuatnya
%  sendiri lalu menarik berkas ini:
%
%      \documentclass[aspectratio=169,11pt]{beamer}
%      % =====================================================================
%  beamer-preamble.tex  --  prasarana dan gaya slide kuliah
%
%  Mata kuliah : Pengantar Kriptografi (MPFA014)
%  Program     : Program Studi Matematika, FMIPA, Universitas Bale Bandung
%
%  Berkas ini TIDAK memuat \documentclass. Setiap berkas minggu memuatnya
%  sendiri lalu menarik berkas ini:
%
%      \documentclass[aspectratio=169,11pt]{beamer}
%      % =====================================================================
%  beamer-preamble.tex  --  prasarana dan gaya slide kuliah
%
%  Mata kuliah : Pengantar Kriptografi (MPFA014)
%  Program     : Program Studi Matematika, FMIPA, Universitas Bale Bandung
%
%  Berkas ini TIDAK memuat \documentclass. Setiap berkas minggu memuatnya
%  sendiri lalu menarik berkas ini:
%
%      \documentclass[aspectratio=169,11pt]{beamer}
%      \input{beamer-preamble}
%      \title{...} \subtitle{...} \setWeekInfo{Minggu n}{Topik}
%
%  Berkas minggu hanya berisi \setWeekInfo dan lingkungan frame.
%  Kompilasi dijalankan DARI DALAM folder ppt/ (lihat compile-ppt.bat),
%  karena gambar buku dirujuk dengan jalur relatif ../buku_ajar/.
% =====================================================================

% ---------------------------------------------------------------- bahasa
\usepackage[utf8]{inputenc}
\usepackage[T1]{fontenc}
\usepackage[indonesian]{babel}
\usepackage{lmodern}
\usepackage{microtype}

% ------------------------------------------------------------- matematika
\usepackage{amsmath,amssymb,mathtools}

% ----------------------------------------------------------------- tabel
\usepackage{booktabs}
\usepackage{array}
\usepackage{tabularx}

% --------------------------------------------------------------- grafik
\usepackage{tikz}
\usetikzlibrary{arrows.meta,positioning,shapes.geometric,calc,fit,backgrounds}

% ---------------------------------------------------------- kode program
\usepackage{listings}
\lstdefinestyle{slidecode}{
    basicstyle=\ttfamily\scriptsize,
    keywordstyle=\color{ubbblue}\bfseries,
    commentstyle=\color{ubbgreen}\itshape,
    stringstyle=\color{ubbred},
    numbers=left,
    numberstyle=\tiny\color{gray},
    numbersep=6pt,
    frame=single,
    rulecolor=\color{ubbblue!40},
    backgroundcolor=\color{ubbblue!3},
    breaklines=true,
    showstringspaces=false,
    tabsize=4,
    columns=fullflexible,
    keepspaces=true,
}
\lstset{style=slidecode}

% =====================================================================
%  WARNA INSTITUSI
% =====================================================================
\definecolor{ubbblue}{RGB}{12, 62, 120}      % warna utama
\definecolor{ubbteal}{RGB}{0, 118, 110}      % aksen kedua
\definecolor{ubbgreen}{RGB}{22, 122, 74}     % untuk blok contoh
\definecolor{ubbred}{RGB}{168, 42, 42}       % untuk blok peringatan
\definecolor{ubbgray}{RGB}{70, 76, 84}       % teks sekunder

% =====================================================================
%  IDENTITAS MATA KULIAH
% =====================================================================
\newcommand{\CourseName}{Pengantar Kriptografi}
\newcommand{\CourseCode}{MPFA014}
\newcommand{\CourseSemester}{Semester Ganjil}
\newcommand{\CourseProgram}{Program Studi Matematika}
\newcommand{\CourseFaculty}{Fakultas Matematika dan Ilmu Pengetahuan Alam}
\newcommand{\CourseInstitution}{Universitas Bale Bandung}
\newcommand{\CourseLecturer}{Cecep Suwanda, S.Si., M.Kom.}

% Label minggu, disetel tiap berkas lewat \setWeekInfo
\newcommand{\WeekLabel}{}
\newcommand{\WeekTopic}{}
\newcommand{\setWeekInfo}[2]{%
    \renewcommand{\WeekLabel}{#1}%
    \renewcommand{\WeekTopic}{#2}%
}

% =====================================================================
%  TEMA: header, footer, warna blok
% =====================================================================
\usetheme{default}
\setbeamertemplate{navigation symbols}{}
\setbeamertemplate{blocks}[rounded][shadow=false]

\setbeamercolor{palette primary}{bg=ubbblue, fg=white}
\setbeamercolor{palette secondary}{bg=ubbgray, fg=white}
\setbeamercolor{palette tertiary}{bg=ubbteal, fg=white}

\setbeamercolor{structure}{fg=ubbblue}
\setbeamercolor{normal text}{fg=black!92}
\setbeamercolor{alerted text}{fg=ubbred}

\setbeamercolor{block title}{bg=ubbblue, fg=white}
\setbeamercolor{block body}{bg=ubbblue!6, fg=black!92}
\setbeamercolor{block title alerted}{bg=ubbred, fg=white}
\setbeamercolor{block body alerted}{bg=ubbred!6, fg=black!92}
\setbeamercolor{block title example}{bg=ubbgreen, fg=white}
\setbeamercolor{block body example}{bg=ubbgreen!6, fg=black!92}

\setbeamercolor{itemize item}{fg=ubbblue}
\setbeamercolor{itemize subitem}{fg=ubbteal}
\setbeamercolor{enumerate item}{fg=ubbblue}

\setbeamerfont{headline}{size=\scriptsize}
\setbeamerfont{footline}{size=\scriptsize}
\setbeamerfont{block title}{size=\normalsize}

% Badan blok rata kiri (bukan rata kanan-kiri): teks slide lebih mudah dibaca
% dan baris terakhir yang pendek tidak lagi memicu "Underfull \hbox".
\addtobeamertemplate{block body begin}{\raggedright}{}
\addtobeamertemplate{block alerted body begin}{\raggedright}{}
\addtobeamertemplate{block example body begin}{\raggedright}{}

% Header: nama mata kuliah dan kode MK
\setbeamertemplate{headline}{%
    \begin{beamercolorbox}[wd=\paperwidth,ht=2.6ex,dp=1.2ex,leftskip=1em,rightskip=1em]{palette primary}%
        \usebeamerfont{headline}\CourseName\ \textperiodcentered\ \CourseCode%
        \hfill \CourseProgram%
    \end{beamercolorbox}%
}

% Footer: minggu dan nomor halaman
\setbeamertemplate{footline}{%
    \begin{beamercolorbox}[wd=\paperwidth,ht=2.4ex,dp=1.2ex,leftskip=1em,rightskip=1em]{palette secondary}%
        \usebeamerfont{footline}\WeekLabel\ \textperiodcentered\ \WeekTopic%
        \hfill \insertframenumber\,/\,\inserttotalframenumber%
    \end{beamercolorbox}%
}

% Halaman judul: judul, subjudul, lalu blok identitas mata kuliah
\setbeamertemplate{title page}{%
    \begin{centering}
        \vskip1.0em
        {\usebeamerfont{title}\usebeamercolor[fg]{title}\inserttitle\par}
        \vskip0.3em
        {\usebeamerfont{subtitle}\usebeamercolor[fg]{subtitle}\insertsubtitle\par}
        \vskip1.2em
        \textcolor{ubbblue}{\rule{0.72\textwidth}{0.8pt}}
        \vskip1.2em
        {\small
            \textbf{\CourseName}\ (\CourseCode)\par
            \vskip0.4em
            \CourseProgram\par
            \CourseFaculty\par
            \CourseInstitution\par
            \vskip0.6em
            \CourseSemester\ \textperiodcentered\ \WeekLabel\par
            \vskip0.8em
            Dosen Pengampu: \CourseLecturer\par
        }
        \vskip1.0em
        \textcolor{ubbblue}{\rule{0.72\textwidth}{0.4pt}}
    \end{centering}
}

% =====================================================================
%  PERINTAH BANTU
% =====================================================================
\newcommand{\R}{\mathbb{R}}
\newcommand{\N}{\mathbb{N}}
\newcommand{\Z}{\mathbb{Z}}
\newcommand{\code}[1]{\texttt{#1}}

% Rujukan gambar dari buku: \bookfig[lebar]{NN}{namaberkas.png}
% Jalur relatif terhadap folder ppt/ (dikompilasi dari sana).
\newcommand{\bookfig}[3][0.80\linewidth]{%
    \includegraphics[width=#1]{../buku_ajar/chapters/#2/figures/#3}%
}

% Keterangan sumber gambar, sepadan dengan \sumber di buku
\newcommand{\figsumber}[1]{%
    \par\vspace{1pt}{\tiny\itshape Sumber: #1\par}%
}

% Penanda blok materi yang sering dipakai di slide
\newcommand{\slidejudul}[1]{{\large\bfseries\textcolor{ubbblue}{#1}}\par\vspace{2pt}}
\newcommand{\kunci}[1]{\textcolor{ubbred}{\textbf{#1}}}

%      \title{...} \subtitle{...} \setWeekInfo{Minggu n}{Topik}
%
%  Berkas minggu hanya berisi \setWeekInfo dan lingkungan frame.
%  Kompilasi dijalankan DARI DALAM folder ppt/ (lihat compile-ppt.bat),
%  karena gambar buku dirujuk dengan jalur relatif ../buku_ajar/.
% =====================================================================

% ---------------------------------------------------------------- bahasa
\usepackage[utf8]{inputenc}
\usepackage[T1]{fontenc}
\usepackage[indonesian]{babel}
\usepackage{lmodern}
\usepackage{microtype}

% ------------------------------------------------------------- matematika
\usepackage{amsmath,amssymb,mathtools}

% ----------------------------------------------------------------- tabel
\usepackage{booktabs}
\usepackage{array}
\usepackage{tabularx}

% --------------------------------------------------------------- grafik
\usepackage{tikz}
\usetikzlibrary{arrows.meta,positioning,shapes.geometric,calc,fit,backgrounds}

% ---------------------------------------------------------- kode program
\usepackage{listings}
\lstdefinestyle{slidecode}{
    basicstyle=\ttfamily\scriptsize,
    keywordstyle=\color{ubbblue}\bfseries,
    commentstyle=\color{ubbgreen}\itshape,
    stringstyle=\color{ubbred},
    numbers=left,
    numberstyle=\tiny\color{gray},
    numbersep=6pt,
    frame=single,
    rulecolor=\color{ubbblue!40},
    backgroundcolor=\color{ubbblue!3},
    breaklines=true,
    showstringspaces=false,
    tabsize=4,
    columns=fullflexible,
    keepspaces=true,
}
\lstset{style=slidecode}

% =====================================================================
%  WARNA INSTITUSI
% =====================================================================
\definecolor{ubbblue}{RGB}{12, 62, 120}      % warna utama
\definecolor{ubbteal}{RGB}{0, 118, 110}      % aksen kedua
\definecolor{ubbgreen}{RGB}{22, 122, 74}     % untuk blok contoh
\definecolor{ubbred}{RGB}{168, 42, 42}       % untuk blok peringatan
\definecolor{ubbgray}{RGB}{70, 76, 84}       % teks sekunder

% =====================================================================
%  IDENTITAS MATA KULIAH
% =====================================================================
\newcommand{\CourseName}{Pengantar Kriptografi}
\newcommand{\CourseCode}{MPFA014}
\newcommand{\CourseSemester}{Semester Ganjil}
\newcommand{\CourseProgram}{Program Studi Matematika}
\newcommand{\CourseFaculty}{Fakultas Matematika dan Ilmu Pengetahuan Alam}
\newcommand{\CourseInstitution}{Universitas Bale Bandung}
\newcommand{\CourseLecturer}{Cecep Suwanda, S.Si., M.Kom.}

% Label minggu, disetel tiap berkas lewat \setWeekInfo
\newcommand{\WeekLabel}{}
\newcommand{\WeekTopic}{}
\newcommand{\setWeekInfo}[2]{%
    \renewcommand{\WeekLabel}{#1}%
    \renewcommand{\WeekTopic}{#2}%
}

% =====================================================================
%  TEMA: header, footer, warna blok
% =====================================================================
\usetheme{default}
\setbeamertemplate{navigation symbols}{}
\setbeamertemplate{blocks}[rounded][shadow=false]

\setbeamercolor{palette primary}{bg=ubbblue, fg=white}
\setbeamercolor{palette secondary}{bg=ubbgray, fg=white}
\setbeamercolor{palette tertiary}{bg=ubbteal, fg=white}

\setbeamercolor{structure}{fg=ubbblue}
\setbeamercolor{normal text}{fg=black!92}
\setbeamercolor{alerted text}{fg=ubbred}

\setbeamercolor{block title}{bg=ubbblue, fg=white}
\setbeamercolor{block body}{bg=ubbblue!6, fg=black!92}
\setbeamercolor{block title alerted}{bg=ubbred, fg=white}
\setbeamercolor{block body alerted}{bg=ubbred!6, fg=black!92}
\setbeamercolor{block title example}{bg=ubbgreen, fg=white}
\setbeamercolor{block body example}{bg=ubbgreen!6, fg=black!92}

\setbeamercolor{itemize item}{fg=ubbblue}
\setbeamercolor{itemize subitem}{fg=ubbteal}
\setbeamercolor{enumerate item}{fg=ubbblue}

\setbeamerfont{headline}{size=\scriptsize}
\setbeamerfont{footline}{size=\scriptsize}
\setbeamerfont{block title}{size=\normalsize}

% Badan blok rata kiri (bukan rata kanan-kiri): teks slide lebih mudah dibaca
% dan baris terakhir yang pendek tidak lagi memicu "Underfull \hbox".
\addtobeamertemplate{block body begin}{\raggedright}{}
\addtobeamertemplate{block alerted body begin}{\raggedright}{}
\addtobeamertemplate{block example body begin}{\raggedright}{}

% Header: nama mata kuliah dan kode MK
\setbeamertemplate{headline}{%
    \begin{beamercolorbox}[wd=\paperwidth,ht=2.6ex,dp=1.2ex,leftskip=1em,rightskip=1em]{palette primary}%
        \usebeamerfont{headline}\CourseName\ \textperiodcentered\ \CourseCode%
        \hfill \CourseProgram%
    \end{beamercolorbox}%
}

% Footer: minggu dan nomor halaman
\setbeamertemplate{footline}{%
    \begin{beamercolorbox}[wd=\paperwidth,ht=2.4ex,dp=1.2ex,leftskip=1em,rightskip=1em]{palette secondary}%
        \usebeamerfont{footline}\WeekLabel\ \textperiodcentered\ \WeekTopic%
        \hfill \insertframenumber\,/\,\inserttotalframenumber%
    \end{beamercolorbox}%
}

% Halaman judul: judul, subjudul, lalu blok identitas mata kuliah
\setbeamertemplate{title page}{%
    \begin{centering}
        \vskip1.0em
        {\usebeamerfont{title}\usebeamercolor[fg]{title}\inserttitle\par}
        \vskip0.3em
        {\usebeamerfont{subtitle}\usebeamercolor[fg]{subtitle}\insertsubtitle\par}
        \vskip1.2em
        \textcolor{ubbblue}{\rule{0.72\textwidth}{0.8pt}}
        \vskip1.2em
        {\small
            \textbf{\CourseName}\ (\CourseCode)\par
            \vskip0.4em
            \CourseProgram\par
            \CourseFaculty\par
            \CourseInstitution\par
            \vskip0.6em
            \CourseSemester\ \textperiodcentered\ \WeekLabel\par
            \vskip0.8em
            Dosen Pengampu: \CourseLecturer\par
        }
        \vskip1.0em
        \textcolor{ubbblue}{\rule{0.72\textwidth}{0.4pt}}
    \end{centering}
}

% =====================================================================
%  PERINTAH BANTU
% =====================================================================
\newcommand{\R}{\mathbb{R}}
\newcommand{\N}{\mathbb{N}}
\newcommand{\Z}{\mathbb{Z}}
\newcommand{\code}[1]{\texttt{#1}}

% Rujukan gambar dari buku: \bookfig[lebar]{NN}{namaberkas.png}
% Jalur relatif terhadap folder ppt/ (dikompilasi dari sana).
\newcommand{\bookfig}[3][0.80\linewidth]{%
    \includegraphics[width=#1]{../buku_ajar/chapters/#2/figures/#3}%
}

% Keterangan sumber gambar, sepadan dengan \sumber di buku
\newcommand{\figsumber}[1]{%
    \par\vspace{1pt}{\tiny\itshape Sumber: #1\par}%
}

% Penanda blok materi yang sering dipakai di slide
\newcommand{\slidejudul}[1]{{\large\bfseries\textcolor{ubbblue}{#1}}\par\vspace{2pt}}
\newcommand{\kunci}[1]{\textcolor{ubbred}{\textbf{#1}}}

%      \title{...} \subtitle{...} \setWeekInfo{Minggu n}{Topik}
%
%  Berkas minggu hanya berisi \setWeekInfo dan lingkungan frame.
%  Kompilasi dijalankan DARI DALAM folder ppt/ (lihat compile-ppt.bat),
%  karena gambar buku dirujuk dengan jalur relatif ../buku_ajar/.
% =====================================================================

% ---------------------------------------------------------------- bahasa
\usepackage[utf8]{inputenc}
\usepackage[T1]{fontenc}
\usepackage[indonesian]{babel}
\usepackage{lmodern}
\usepackage{microtype}

% ------------------------------------------------------------- matematika
\usepackage{amsmath,amssymb,mathtools}

% ----------------------------------------------------------------- tabel
\usepackage{booktabs}
\usepackage{array}
\usepackage{tabularx}

% --------------------------------------------------------------- grafik
\usepackage{tikz}
\usetikzlibrary{arrows.meta,positioning,shapes.geometric,calc,fit,backgrounds}

% ---------------------------------------------------------- kode program
\usepackage{listings}
\lstdefinestyle{slidecode}{
    basicstyle=\ttfamily\scriptsize,
    keywordstyle=\color{ubbblue}\bfseries,
    commentstyle=\color{ubbgreen}\itshape,
    stringstyle=\color{ubbred},
    numbers=left,
    numberstyle=\tiny\color{gray},
    numbersep=6pt,
    frame=single,
    rulecolor=\color{ubbblue!40},
    backgroundcolor=\color{ubbblue!3},
    breaklines=true,
    showstringspaces=false,
    tabsize=4,
    columns=fullflexible,
    keepspaces=true,
}
\lstset{style=slidecode}

% =====================================================================
%  WARNA INSTITUSI
% =====================================================================
\definecolor{ubbblue}{RGB}{12, 62, 120}      % warna utama
\definecolor{ubbteal}{RGB}{0, 118, 110}      % aksen kedua
\definecolor{ubbgreen}{RGB}{22, 122, 74}     % untuk blok contoh
\definecolor{ubbred}{RGB}{168, 42, 42}       % untuk blok peringatan
\definecolor{ubbgray}{RGB}{70, 76, 84}       % teks sekunder

% =====================================================================
%  IDENTITAS MATA KULIAH
% =====================================================================
\newcommand{\CourseName}{Pengantar Kriptografi}
\newcommand{\CourseCode}{MPFA014}
\newcommand{\CourseSemester}{Semester Ganjil}
\newcommand{\CourseProgram}{Program Studi Matematika}
\newcommand{\CourseFaculty}{Fakultas Matematika dan Ilmu Pengetahuan Alam}
\newcommand{\CourseInstitution}{Universitas Bale Bandung}
\newcommand{\CourseLecturer}{Cecep Suwanda, S.Si., M.Kom.}

% Label minggu, disetel tiap berkas lewat \setWeekInfo
\newcommand{\WeekLabel}{}
\newcommand{\WeekTopic}{}
\newcommand{\setWeekInfo}[2]{%
    \renewcommand{\WeekLabel}{#1}%
    \renewcommand{\WeekTopic}{#2}%
}

% =====================================================================
%  TEMA: header, footer, warna blok
% =====================================================================
\usetheme{default}
\setbeamertemplate{navigation symbols}{}
\setbeamertemplate{blocks}[rounded][shadow=false]

\setbeamercolor{palette primary}{bg=ubbblue, fg=white}
\setbeamercolor{palette secondary}{bg=ubbgray, fg=white}
\setbeamercolor{palette tertiary}{bg=ubbteal, fg=white}

\setbeamercolor{structure}{fg=ubbblue}
\setbeamercolor{normal text}{fg=black!92}
\setbeamercolor{alerted text}{fg=ubbred}

\setbeamercolor{block title}{bg=ubbblue, fg=white}
\setbeamercolor{block body}{bg=ubbblue!6, fg=black!92}
\setbeamercolor{block title alerted}{bg=ubbred, fg=white}
\setbeamercolor{block body alerted}{bg=ubbred!6, fg=black!92}
\setbeamercolor{block title example}{bg=ubbgreen, fg=white}
\setbeamercolor{block body example}{bg=ubbgreen!6, fg=black!92}

\setbeamercolor{itemize item}{fg=ubbblue}
\setbeamercolor{itemize subitem}{fg=ubbteal}
\setbeamercolor{enumerate item}{fg=ubbblue}

\setbeamerfont{headline}{size=\scriptsize}
\setbeamerfont{footline}{size=\scriptsize}
\setbeamerfont{block title}{size=\normalsize}

% Badan blok rata kiri (bukan rata kanan-kiri): teks slide lebih mudah dibaca
% dan baris terakhir yang pendek tidak lagi memicu "Underfull \hbox".
\addtobeamertemplate{block body begin}{\raggedright}{}
\addtobeamertemplate{block alerted body begin}{\raggedright}{}
\addtobeamertemplate{block example body begin}{\raggedright}{}

% Header: nama mata kuliah dan kode MK
\setbeamertemplate{headline}{%
    \begin{beamercolorbox}[wd=\paperwidth,ht=2.6ex,dp=1.2ex,leftskip=1em,rightskip=1em]{palette primary}%
        \usebeamerfont{headline}\CourseName\ \textperiodcentered\ \CourseCode%
        \hfill \CourseProgram%
    \end{beamercolorbox}%
}

% Footer: minggu dan nomor halaman
\setbeamertemplate{footline}{%
    \begin{beamercolorbox}[wd=\paperwidth,ht=2.4ex,dp=1.2ex,leftskip=1em,rightskip=1em]{palette secondary}%
        \usebeamerfont{footline}\WeekLabel\ \textperiodcentered\ \WeekTopic%
        \hfill \insertframenumber\,/\,\inserttotalframenumber%
    \end{beamercolorbox}%
}

% Halaman judul: judul, subjudul, lalu blok identitas mata kuliah
\setbeamertemplate{title page}{%
    \begin{centering}
        \vskip1.0em
        {\usebeamerfont{title}\usebeamercolor[fg]{title}\inserttitle\par}
        \vskip0.3em
        {\usebeamerfont{subtitle}\usebeamercolor[fg]{subtitle}\insertsubtitle\par}
        \vskip1.2em
        \textcolor{ubbblue}{\rule{0.72\textwidth}{0.8pt}}
        \vskip1.2em
        {\small
            \textbf{\CourseName}\ (\CourseCode)\par
            \vskip0.4em
            \CourseProgram\par
            \CourseFaculty\par
            \CourseInstitution\par
            \vskip0.6em
            \CourseSemester\ \textperiodcentered\ \WeekLabel\par
            \vskip0.8em
            Dosen Pengampu: \CourseLecturer\par
        }
        \vskip1.0em
        \textcolor{ubbblue}{\rule{0.72\textwidth}{0.4pt}}
    \end{centering}
}

% =====================================================================
%  PERINTAH BANTU
% =====================================================================
\newcommand{\R}{\mathbb{R}}
\newcommand{\N}{\mathbb{N}}
\newcommand{\Z}{\mathbb{Z}}
\newcommand{\code}[1]{\texttt{#1}}

% Rujukan gambar dari buku: \bookfig[lebar]{NN}{namaberkas.png}
% Jalur relatif terhadap folder ppt/ (dikompilasi dari sana).
\newcommand{\bookfig}[3][0.80\linewidth]{%
    \includegraphics[width=#1]{../buku_ajar/chapters/#2/figures/#3}%
}

% Keterangan sumber gambar, sepadan dengan \sumber di buku
\newcommand{\figsumber}[1]{%
    \par\vspace{1pt}{\tiny\itshape Sumber: #1\par}%
}

% Penanda blok materi yang sering dipakai di slide
\newcommand{\slidejudul}[1]{{\large\bfseries\textcolor{ubbblue}{#1}}\par\vspace{2pt}}
\newcommand{\kunci}[1]{\textcolor{ubbred}{\textbf{#1}}}


\title{One-Time Pad: Rahasia yang Terbukti Sempurna}
\subtitle{OTP Vernam, tiga syarat kerahasiaan sempurna, dan kelemahan praktisnya}
\setWeekInfo{Minggu 6}{One-time pad}

\begin{document}

% ----------------------------------------------------------- halaman judul
\begin{frame}[plain]
    \titlepage
\end{frame}

% ------------------------------------------------------------------ peta
\begin{frame}{Peta Materi Pertemuan Ini}
    \footnotesize
    \begin{columns}[T,onlytextwidth]
        \begin{column}{0.54\textwidth}
            \slidejudul{Enam bagian}
            \begin{enumerate}
                \item Batas cipher klasik dan tiga syaratnya
                \item Prinsip kerja OTP dan aturan modular
                \item Contoh hitung sepuluh huruf
                \item Tiga syarat kerahasiaan sempurna
                \item Kerahasiaan sempurna dan kelemahan praktis
                \item Sejarah, rangkuman, dan latihan
            \end{enumerate}
        \end{column}
        \begin{column}{0.42\textwidth}
            \begin{block}{Referensi}
                \footnotesize
                Bab 4 buku ajar; Munir, \textit{IF4020 Kriptografi}, ITB,
                2026 (modul 05).
            \end{block}
        \end{column}
    \end{columns}
    \vspace{0.15cm}
    \begin{exampleblock}{Benang merah}
        OTP adalah satu-satunya cipher yang keamanannya dapat dibuktikan
        secara matematis --- dan harganya paling mahal.
    \end{exampleblock}
\end{frame}

% ---------------------------------------------------- capaian pembelajaran
\begin{frame}{Capaian Pembelajaran}
    \begin{block}{Sub-CPMK 1.2 --- One-time pad dan stream cipher modern}
        \footnotesize
        Mengenal one-time pad serta stream cipher modern, yaitu representasi
        bit, operasi XOR, keystream generator, FSR, RC4, A5, Trivium,
        Salsa20, dan ChaCha20.
    \end{block}
    \vspace{0.1cm}
    {\scriptsize\itshape Pertemuan ini menuntaskan bagian one-time pad;
    bagian stream cipher modern menjadi bahan Minggu 7.}
    \vspace{0.15cm}
    \begin{exampleblock}{Setelah pertemuan ini, mahasiswa mampu}
        \footnotesize
        \begin{itemize}
            \item menjelaskan mengapa cipher klasik dapat dipecahkan dan
                  apa yang menutup celah itu;
            \item menyebutkan tiga syarat cipher yang tidak dapat
                  dipecahkan;
            \item menerapkan aturan $C=(P+K)\bmod 26$ dan
                  $P=(C-K)\bmod 26$ pada contoh buku ajar;
            \item menjelaskan makna kerahasiaan sempurna menurut Shannon;
            \item mengidentifikasi kelemahan praktis OTP dan alasan OTP
                  jarang dipakai.
        \end{itemize}
    \end{exampleblock}
\end{frame}

% =====================================================================
\section{Menuju Cipher yang Tidak Dapat Dipecahkan}
% =====================================================================

\begin{frame}{Klaim yang Tidak Pernah Terbukti}
    \begin{columns}[T,onlytextwidth]
        \begin{column}{0.56\textwidth}
            \footnotesize
            Setiap perancang cipher berharap karyanya \emph{unbreakable}.
            Kenyataannya, hampir semua cipher yang pernah dibuat
            \emph{breakable}:
            \begin{itemize}
                \item Caesar, Vigen\`ere, Playfair, dan Affine;
                \item Enigma dan Hill.
            \end{itemize}
            Semuanya sudah kedaluwarsa karena berhasil dipecahkan.
        \end{column}
        \begin{column}{0.40\textwidth}
            \begin{alertblock}{Pertanyaan}
                \footnotesize
                Adakah cipher yang benar-benar tidak dapat dipecahkan?
            \end{alertblock}
        \end{column}
    \end{columns}
    \vspace{0.25cm}
    \begin{block}{Jawaban singkat}
        Ada, tetapi hanya bila tiga syarat dipenuhi sekaligus.
    \end{block}
\end{frame}

\begin{frame}{Tiga Syarat Cipher yang Tidak Dapat Dipecahkan}
    \footnotesize
    \begin{enumerate}
        \item \textbf{Kunci acak sejati.} Nilainya tidak dapat diprediksi
              dan pembangkitannya tidak dapat diulang.
        \item \textbf{Panjang kunci sama dengan panjang plainteks.}
              Tidak ada huruf kunci yang dipakai untuk dua huruf plainteks.
        \item \textbf{Kunci hanya dipakai sekali.} Setelah satu pesan,
              kunci itu tidak boleh muncul lagi.
    \end{enumerate}
    \vspace{0.25cm}
    \begin{alertblock}{Ketiganya satu paket}
        Melanggar satu syarat saja sudah cukup untuk membatalkan jaminan
        kerahasiaan. Dua syarat yang lain tidak dapat menutupinya.
    \end{alertblock}
\end{frame}

\begin{frame}{Acak Sejati dan Acak Semu}
    \begin{columns}[T,onlytextwidth]
        \begin{column}{0.48\textwidth}
            \begin{block}{Acak sejati}
                \footnotesize
                \begin{itemize}
                    \item nilainya tidak dapat diprediksi;
                    \item pembangkitannya tidak dapat diulang.
                \end{itemize}
            \end{block}
        \end{column}
        \begin{column}{0.48\textwidth}
            \begin{block}{Acak semu}
                \footnotesize
                \begin{itemize}
                    \item dihasilkan oleh prosedur tertentu;
                    \item dapat diulang bila umpan (\textit{seed}) sama.
                \end{itemize}
            \end{block}
        \end{column}
    \end{columns}
    \vspace{0.15cm}
    \begin{exampleblock}{Akibat dua syarat pertama}
        Plainteks yang sama tidak selalu menghasilkan cipherteks yang sama.
        Karena itu OTP menuntut acak sejati, bukan acak semu.
    \end{exampleblock}
\end{frame}

% =====================================================================
\section{Prinsip Kerja One-Time Pad}
% =====================================================================

\begin{frame}{Vernam 1917 dan Shannon 1949}
    \begin{columns}[T,onlytextwidth]
        \begin{column}{0.58\textwidth}
            \footnotesize
            \begin{itemize}
                \item Gilbert Vernam menemukan OTP pada tahun 1917.
                \item Claude Shannon membuktikan keamanannya secara
                      matematis pada tahun 1949.
                \item OTP menjawab kelemahan Vigen\`ere: kunci tidak lagi
                      diulang secara periodik.
            \end{itemize}
            \vspace{0.1cm}
            \begin{exampleblock}{Yang membuat OTP berbeda}
                Panjang kunci menyamai panjang pesan dan seluruh isinya acak
                sejati, sehingga tidak ada pola yang dapat ditelusuri.
            \end{exampleblock}
        \end{column}
        \begin{column}{0.38\textwidth}
            \centering
            \bookfig[0.62\linewidth]{04}{gilbert-vernam.png}
            \figsumber{Munir, \textit{IF4020 Kriptografi}, ITB.}
        \end{column}
    \end{columns}
\end{frame}

\begin{frame}{Pad Kertas yang Hanya Dipakai Sekali}
    \footnotesize
    Nama \emph{one-time pad} mengacu pada bloknot kertas berisi deretan
    huruf kunci yang dibangkitkan secara acak.
    \begin{itemize}
        \item Pengirim dan penerima memegang salinan pad yang sama.
        \item Satu pad dipakai untuk satu pesan saja.
        \item Setelah dipakai, pad dimusnahkan agar tidak terpakai lagi.
    \end{itemize}
    \vspace{0.2cm}
    \begin{alertblock}{Satu pad, satu pesan}
        Aturan ini bukan sekadar anjuran. Memakai ulang pad mengubah OTP
        menjadi \emph{two-time pad}, yang tidak lagi aman.
    \end{alertblock}
\end{frame}

\begin{frame}{Aturan Enkripsi dan Dekripsi}
    \footnotesize
    Setiap huruf dipetakan ke bilangan $0$ sampai $25$. Aturannya sama
    dengan Vigen\`ere, tetapi tanpa pengulangan kunci:
    \[
        C = (P + K) \bmod 26, \qquad
        P = (C - K) \bmod 26 .
    \]
    \begin{itemize}
        \item Panjang kunci sama dengan panjang plainteks.
        \item Setiap huruf plainteks berpasangan dengan satu huruf kunci
              yang khusus untuknya.
        \item Tidak ada kunci yang dipakai dua kali di dalam satu pesan.
    \end{itemize}
    \vspace{0.1cm}
    \begin{exampleblock}{Akibatnya}
        Metode Kasiski mencari jarak pengulangan kunci. Pada OTP jarak itu
        tidak pernah ada, sehingga metodenya kehilangan pijakan.
    \end{exampleblock}
\end{frame}

\begin{frame}{Contoh Satu Huruf}
    \begin{columns}[T,onlytextwidth]
        \begin{column}{0.52\textwidth}
            \footnotesize
            Ambil $P=\code{o}$ dan $K=\code{t}$. Dalam bilangan,
            $P=14$ dan $K=19$.
            \begin{align*}
                C &= (14+19) \bmod 26 \\
                  &= 33 \bmod 26 = 7 = \code{H}.
            \end{align*}
            Dekripsi mengembalikan huruf semula:
            \begin{align*}
                P &= (7-19) \bmod 26 \\
                  &= -12 \bmod 26 = 14 = \code{o}.
            \end{align*}
        \end{column}
        \begin{column}{0.44\textwidth}
            \centering
            \bookfig[0.98\linewidth]{04}{contoh-one-time-pad.png}
            \figsumber{Munir, \textit{IF4020 Kriptografi}, ITB.}
        \end{column}
    \end{columns}
\end{frame}

% =====================================================================
\section{Contoh Perhitungan dari Buku Ajar}
% =====================================================================

\begin{frame}{Data Contoh}
    \footnotesize
    Buku ajar memakai plainteks dan kunci yang sama-sama sepuluh huruf.
    Tidak ada huruf kunci yang berulang menurut posisi:
    \begin{center}
    \small
    \begin{tabular}{@{}l l@{}}
        \toprule
        Plainteks  & \code{onetimepad} \\
        Kunci      & \code{tbfrgfarhm} \\
        \midrule
        Cipherteks & \code{HOJKOREGHP} \\
        \bottomrule
    \end{tabular}
    \end{center}
    \vspace{0.1cm}
    \begin{exampleblock}{Huruf pertama}
        $(\code{o}+\code{t}) = (14+19) = 33$, lalu
        $33 \bmod 26 = 7 = \code{H}$.
    \end{exampleblock}
    \begin{exampleblock}{Huruf kedua}
        $(\code{n}+\code{b}) = (13+1) = 14$, lalu
        $14 \bmod 26 = 14 = \code{O}$.
    \end{exampleblock}
\end{frame}

\begin{frame}{Sepuluh Langkah Enkripsi}
    \scriptsize
    Setiap baris adalah satu huruf. Operasinya
    $c_i = (p_i + k_i) \bmod 26$, dengan $p_i$ dan $k_i$ dalam bilangan
    $0$ sampai $25$.
    \begin{center}
    \setlength{\tabcolsep}{6pt}
    \begin{tabular}{@{}c c c c c c@{}}
        \toprule
        $i$ & $p_i$ & $k_i$ & $p_i+k_i$ & $\bmod\,26$ & $c_i$ \\
        \midrule
        1  & \code{o} (14) & \code{t} (19) & 33 & 7  & \code{H} \\
        2  & \code{n} (13) & \code{b} (1)  & 14 & 14 & \code{O} \\
        3  & \code{e} (4)  & \code{f} (5)  & 9  & 9  & \code{J} \\
        4  & \code{t} (19) & \code{r} (17) & 36 & 10 & \code{K} \\
        5  & \code{i} (8)  & \code{g} (6)  & 14 & 14 & \code{O} \\
        6  & \code{m} (12) & \code{f} (5)  & 17 & 17 & \code{R} \\
        7  & \code{e} (4)  & \code{a} (0)  & 4  & 4  & \code{E} \\
        8  & \code{p} (15) & \code{r} (17) & 32 & 6  & \code{G} \\
        9  & \code{a} (0)  & \code{h} (7)  & 7  & 7  & \code{H} \\
        10 & \code{d} (3)  & \code{m} (12) & 15 & 15 & \code{P} \\
        \bottomrule
    \end{tabular}
    \end{center}
    \vspace{0.1cm}
    Baris pertama memperlihatkan urutan lengkapnya:
    $14+19=33$, $33 \bmod 26=7$, dan $7$ adalah huruf \code{H}.
    Hasil kesepuluh baris itu adalah \code{HOJKOREGHP}.
\end{frame}

\begin{frame}{Pemeriksaan Dekripsi}
    \footnotesize
    Dekripsi memakai pengurangan, $p_i=(c_i-k_i)\bmod 26$. Tiga huruf
    pertama:
    \begin{align*}
        (\code{H}-\code{t}) &= (7-19) \bmod 26 = 14 = \code{o}, \\
        (\code{O}-\code{b}) &= (14-1) \bmod 26 = 13 = \code{n}, \\
        (\code{J}-\code{f}) &= (9-5) \bmod 26 = 4 = \code{e}.
    \end{align*}
    Bila seluruh baris dikerjakan, hasilnya kembali menjadi
    \code{onetimepad}.
    \vspace{0.1cm}
    \begin{exampleblock}{Yang perlu diperhatikan}
        Sepuluh huruf memakai sepuluh huruf kunci yang berbeda. Tidak ada
        satu operasi pun yang berlaku untuk dua posisi sekaligus.
    \end{exampleblock}
\end{frame}

% =====================================================================
\section{Tiga Syarat Kerahasiaan Sempurna}
% =====================================================================

\begin{frame}{Syarat 1: Kunci Acak Sejati}
    \footnotesize
    Kunci OTP harus seluruhnya acak sejati dan sepanjang pesan.
    \begin{itemize}
        \item Kunci yang diambil dari novel, buku, atau berita tidak acak.
        \item Teks semacam itu dapat ditebak dan mengandung pola bahasa.
        \item Akibatnya cipher tersebut bukan OTP dan tidak memiliki
              kerahasiaan sempurna.
    \end{itemize}
    \vspace{0.2cm}
    \begin{alertblock}{Panjang saja tidak cukup}
        Kunci sepanjang pesan tetap lemah bila isinya dapat dibaca atau
        diperkirakan penyerang.
    \end{alertblock}
\end{frame}

\begin{frame}{Syarat 2: Panjang Kunci Sama dengan Plainteks}
    \footnotesize
    \begin{itemize}
        \item Setiap huruf plainteks memperoleh huruf kunci sendiri.
        \item Tidak ada pengulangan kunci secara periodik.
        \item Pola jarak yang biasa ditelusuri Kasiski tidak terbentuk.
    \end{itemize}
    \vspace{0.2cm}
    \begin{exampleblock}{Bandingkan dengan Vigen\`ere}
        Pada Vigen\`ere kunci pendek diulang, sehingga cipherteks menyimpan
        jejak periodisitas. Pada OTP jejak itu dihapus oleh panjang kunci.
    \end{exampleblock}
\end{frame}

\begin{frame}{Syarat 3: Kunci Hanya Dipakai Sekali}
    \footnotesize
    Bila kunci yang sama dipakai untuk dua pesan, cipher itu bukan lagi
    \emph{one-time pad}, melainkan \emph{two-time pad}.
    \[
        C_1 - C_2 = (P_1 + K) - (P_2 + K) = P_1 - P_2 \pmod{26}.
    \]
    Kunci saling menghapus, sehingga yang tersisa hanya selisih kedua
    plainteks. Penyerang tidak perlu mengetahui kunci untuk memperoleh
    hubungan itu.
    \vspace{0.1cm}
    \begin{alertblock}{Kesimpulan}
        Pemakaian ulang kunci membuat cipher tersebut tidak aman, sepanjang
        apa pun pesannya.
    \end{alertblock}
\end{frame}

\begin{frame}{Akibatnya: Pemetaan Satu-ke-Satu}
    \footnotesize
    Untuk sembarang plainteks $P$ dan cipherteks $C$, hanya ada satu kunci
    yang memetakan keduanya:
    \[
        K = (C - P) \bmod 26 .
    \]
    \begin{itemize}
        \item Kriptanalis dapat mendekripsi $C$ dengan beberapa kunci
              berbeda.
        \item Hasilnya beberapa plainteks yang semuanya bermakna.
        \item Ia tidak dapat menentukan plainteks mana yang benar.
    \end{itemize}
    \vspace{0.15cm}
    \begin{exampleblock}{Inti argumennya}
        Cipherteks tidak menyempitkan ruang kemungkinan plainteks. Semua
        tebakan tetap sama masuk akal.
    \end{exampleblock}
\end{frame}

\begin{frame}{Satu Cipherteks, Dua Plainteks Bermakna}
    \footnotesize
    Cipherteks \code{HOJKOREGHP} dapat dijelaskan oleh lebih dari satu
    kunci. Dua di antaranya:
    \begin{center}
    \small
    \setlength{\tabcolsep}{6pt}
    \begin{tabular}{@{}l l l@{}}
        \toprule
        \textbf{Cipherteks} & \textbf{Kunci} & \textbf{Plainteks} \\
        \midrule
        \code{HOJKOREGHP} & \code{POYYAEAABX} & \code{SALMONEGGS} \\
        \code{HOJKOREGHP} & \code{BXFGBMWCWM} & \code{GREENFIELD} \\
        \bottomrule
    \end{tabular}
    \end{center}
    \vspace{0.1cm}
    Kedua kunci dihitung dengan $K=(C-P)\bmod 26$. Keduanya sah dan
    keduanya menghasilkan kalimat yang bermakna.
    \vspace{0.1cm}
    \begin{alertblock}{Yang tidak dapat dilakukan penyerang}
        Memilih di antara keduanya hanya dari cipherteks. Ia membutuhkan
        informasi lain di luar cipherteks.
    \end{alertblock}
\end{frame}

% =====================================================================
\section{Kerahasiaan Sempurna}
% =====================================================================

\begin{frame}{Definisi Shannon}
    \begin{columns}[T,onlytextwidth]
        \begin{column}{0.58\textwidth}
            \footnotesize
            Pada tahun 1949 Claude Shannon, dari Laboratorium Bell,
            membuktikan bahwa OTP memiliki \emph{perfect secrecy}.
            \begin{itemize}
                \item Buktinya dimuat dalam \emph{Communication Theory of
                      Secrecy Systems}.
                \item Setelah cipherteks diperoleh, informasi tentang
                      plainteks tidak bertambah.
            \end{itemize}
            \[
                P(P=p \mid C=c) = P(P=p).
            \]
            Probabilitas plainteks sama sebelum dan sesudah cipherteks
            diamati.
        \end{column}
        \begin{column}{0.38\textwidth}
            \centering
            \bookfig[0.62\linewidth]{04}{claude-shannon.png}
            \figsumber{Munir, \textit{IF4020 Kriptografi}, ITB.}
        \end{column}
    \end{columns}
\end{frame}

\begin{frame}{Mengapa OTP Memenuhi Definisi Itu}
    \footnotesize
    Misalkan kunci dipilih uniform, sehingga $P(K=k)=1/26$ untuk
    $k=0,1,\dots,25$.
    \begin{itemize}
        \item Untuk setiap plainteks $p$, pemetaan $C=(P+K)\bmod 26$
              adalah satu-ke-satu dari semua kunci ke semua cipherteks.
        \item Karena itu $P(C=c \mid P=p)=1/26$ untuk setiap $c$.
        \item Setiap huruf cipherteks sama mungkinnya, tidak bergantung
              pada $p$.
    \end{itemize}
    \vspace{0.15cm}
    \begin{exampleblock}{Maknanya}
        Cipherteks tidak memberi petunjuk untuk mempersempit tebakan
        plainteks. Inilah yang membuat OTP disebut sempurna secara teori.
    \end{exampleblock}
\end{frame}

% =====================================================================
\section{Kelemahan Praktis dan Warisan Sejarah}
% =====================================================================

\begin{frame}{Kelemahan 1: Kunci Sepanjang Pesan}
    \footnotesize
    \begin{itemize}
        \item Ukuran kunci tumbuh seiring panjang pesan.
        \item Pesan yang sangat panjang menuntut kunci yang sama panjang.
        \item Membangkitkan miliaran karakter yang benar-benar acak
              membutuhkan komputasi berat.
        \item Komunikasi yang berlangsung terus-menerus menuntut pasokan
              kunci tanpa henti.
    \end{itemize}
    \vspace{0.2cm}
    \begin{alertblock}{Sifat yang melekat}
        Ini bukan cacat implementasi, melainkan harga langsung dari syarat
        panjang kunci sama dengan plainteks.
    \end{alertblock}
\end{frame}

\begin{frame}{Kelemahan 2: Kunci Harus Dibangkitkan Acak Sejati}
    \footnotesize
    \begin{itemize}
        \item Pengirim dan penerima tidak dapat membangkitkan kunci acak
              yang sama secara bersamaan.
        \item Keduanya harus menyepakati kunci lebih dahulu.
        \item Kunci itu harus sepanjang pesan dan hanya boleh dipakai
              sekali.
    \end{itemize}
    \vspace{0.2cm}
    \begin{exampleblock}{Akibatnya}
        Kesepakatan kunci menjadi pekerjaan tersendiri yang harus
        diselesaikan sebelum pesan pertama dikirim.
    \end{exampleblock}
\end{frame}

\begin{frame}{Distribusi Kunci Butuh Saluran Kedua}
    \footnotesize
    OTP hanya dapat dipakai bila tersedia saluran komunikasi kedua yang
    sangat aman untuk mengirim kunci.
    \begin{itemize}
        \item Saluran itu tidak boleh sama dengan saluran pesan.
        \item Umumnya lambat dan mahal, misalnya kurir terpercaya lewat
              jalur darat.
        \item Alternatifnya, kunci dikirim pada saluran yang sama tetapi
              pada waktu yang berbeda.
    \end{itemize}
    \vspace{0.2cm}
    \begin{alertblock}{Pertukaran yang tidak mudah}
        Keamanan sempurna di saluran pesan dibayar dengan kebutuhan
        saluran rahasia untuk kuncinya.
    \end{alertblock}
\end{frame}

\begin{frame}{Tidak Cocok untuk Komunikasi Dinamis}
    \footnotesize
    Komunikasi modern yang terus berubah sulit dilayani OTP:
    \begin{itemize}
        \item WhatsApp, HTTPS, VPN, dan \textit{cloud storage};
        \item email dan komunikasi IoT.
    \end{itemize}
    Pada komunikasi semacam itu data terus bertambah, perangkat dapat
    berganti, pengguna dapat \textit{offline}, dan ukuran pesan
    berbeda-beda.
    \vspace{0.1cm}
    \begin{alertblock}{Beban manajemen kunci}
        Setiap tambahan data memerlukan tambahan kunci acak yang belum
        pernah dipakai. Manajemen kunci menjadi sangat berat.
    \end{alertblock}
\end{frame}

\begin{frame}{Syal Agen Rahasia Perang Dunia II}
    \begin{columns}[T,onlytextwidth]
        \begin{column}{0.46\textwidth}
            \footnotesize
            OTP pernah dipakai oleh mata-mata selama Perang Dunia II.
            \begin{itemize}
                \item Kunci dicetak pada kain, bukan kertas.
                \item Kain mudah disembunyikan.
                \item Kain juga mudah dimusnahkan setelah dipakai.
            \end{itemize}
            Pejabat eksekutif operasi khusus Inggris memperlihatkan syal
            berisi kunci OTP sebagai contohnya.
        \end{column}
        \begin{column}{0.50\textwidth}
            \centering
            \bookfig[0.78\linewidth]{04}{otp-syal.png}
            \figsumber{Munir, \textit{IF4020 Kriptografi}, ITB.}
        \end{column}
    \end{columns}
\end{frame}

\begin{frame}{Hotline Moskow--Washington 1963}
    \begin{columns}[T,onlytextwidth]
        \begin{column}{0.42\textwidth}
            \centering
            \bookfig[0.55\linewidth]{04}{otp-hotline.png}
            \figsumber{Munir, \textit{IF4020 Kriptografi}, ITB.}
        \end{column}
        \begin{column}{0.54\textwidth}
            \footnotesize
            Saluran langsung antara pemimpin Amerika Serikat dan Uni Soviet
            dibangun setelah Krisis Rudal Kuba tahun 1963.
            \begin{itemize}
                \item Peralatannya OTP bernama ITT Intelex Teletype L015.
                \item Pita kunci \emph{one-time pad} disimpan dan
                      diterbangkan antara Washington DC dan Moskow.
            \end{itemize}
            Kunci tidak pernah dikirim melalui saluran pesan itu sendiri.
        \end{column}
    \end{columns}
\end{frame}

\begin{frame}{Penilaian Bellovin}
    \footnotesize
    Steve Bellovin merangkum ketegangan antara teori dan praktik:
    \begin{columns}[T,onlytextwidth]
        \begin{column}{0.48\textwidth}
            \begin{block}{OTP}
                Tidak dapat dipecahkan secara teori, tetapi rapuh dalam
                praktik karena beban kuncinya.
            \end{block}
        \end{column}
        \begin{column}{0.48\textwidth}
            \begin{block}{Cipher konvensional}
                Dapat dipecahkan secara teori, tetapi kuat dalam praktik
                karena kuncinya pendek dan mudah dikelola.
            \end{block}
        \end{column}
    \end{columns}
    \vspace{0.2cm}
    \begin{exampleblock}{Pelajaran}
        Keamanan teoretis tidak otomatis menjadi keamanan operasional.
        Pilihan cipher selalu menimbang keduanya.
    \end{exampleblock}
\end{frame}

% =====================================================================
\section{Rangkuman dan Latihan}
% =====================================================================

\begin{frame}{Rangkuman (1/2)}
    \footnotesize
    \begin{itemize}
        \item Cipher klasik berguguran karena menyimpan pola. OTP menutup
              pola itu dengan kunci acak sejati sepanjang pesan.
        \item Tiga syaratnya: acak sejati, panjang sama dengan plainteks,
              dan hanya dipakai sekali.
        \item Enkripsi memakai $C=(P+K)\bmod 26$; dekripsi memakai
              $P=(C-K)\bmod 26$.
        \item Contoh buku ajar: \code{onetimepad} dengan kunci
              \code{tbfrgfarhm} menghasilkan \code{HOJKOREGHP}.
    \end{itemize}
\end{frame}

\begin{frame}{Rangkuman (2/2)}
    \footnotesize
    \begin{itemize}
        \item Kerahasiaan sempurna berarti
              $P(P=p \mid C=c)=P(P=p)$: cipherteks tidak menambah
              informasi tentang plainteks.
        \item Pemetaan satu-ke-satu membuat banyak plainteks bermakna
              mungkin; kriptanalis tidak dapat memilih yang benar.
        \item Kelemahan praktisnya: kunci sepanjang pesan, pembangkitan
              acak sejati yang berat, distribusi lewat saluran kedua, dan
              ketidakcocokan dengan komunikasi dinamis.
        \item OTP terbukti istimewa secara teori, tetapi bukan pilihan
              sehari-hari.
    \end{itemize}
\end{frame}

\begin{frame}{Latihan}
    \footnotesize
    \begin{enumerate}
        \item Dengan $P=\code{onetimepad}$ dan $K=\code{tbfrgfarhm}$,
              hitunglah cipherteksnya huruf demi huruf.
        \item Dekripsikan \code{HOJKOREGHP} dengan kunci
              \code{tbfrgfarhm}, lalu bandingkan hasilnya dengan plainteks
              semula.
        \item Mengapa kunci yang diambil dari sebuah novel tidak
              menghasilkan kerahasiaan sempurna?
        \item Jelaskan bagaimana dua kunci berbeda dapat menghasilkan dua
              plainteks bermakna dari cipherteks yang sama.
        \item Sebutkan tiga kelemahan praktis OTP dan tentukan mana yang
              paling menghambat komunikasi modern.
    \end{enumerate}
\end{frame}

\begin{frame}{Penutup}
    \begin{block}{Tiga hal yang harus dibawa dari pertemuan ini}
        \begin{enumerate}
            \item OTP adalah cipher yang keamanannya terbukti, dengan syarat
                  kunci acak sejati, sepanjang pesan, dan sekali pakai.
            \item Kerahasiaan sempurna berarti cipherteks tidak memberi
                  informasi tambahan tentang plainteks.
            \item Kelemahan OTP bersifat praktis: kunci besar, distribusi
                  sulit, dan pemakaian ulang menghancurkan keamanannya.
        \end{enumerate}
    \end{block}
    \vspace{0.3cm}
    \begin{exampleblock}{Pertemuan berikutnya --- Minggu 7}
        \textbf{Minggu 7 --- Kriptografi modern dan stream cipher}
    \end{exampleblock}
\end{frame}

\end{document}
